\subsection{Polynomial Functional Calculus}

\begin{theorem} \label{thm:polynomial-functional-calculus-diagonalizable-matrix} \hlt{Polynomial Functional Calculus for Diagonalizable Matrices}\\
Let $A\in\operatorname{Mat}(n,n,\mathbb{F})$ be diagonalizable and write $A=S\Lambda S^{-1}$, where $\Lambda=\operatorname{diag}(\lambda_1,\ldots,\lambda_n)$.\\
If $p\in\mathbb{F}[z]$, then $p(A)=S\operatorname{diag}(p(\lambda_1),\ldots,p(\lambda_n))S^{-1}$.
\begin{proof}
For every $j\geq0$, $A^j=S\Lambda^jS^{-1}$.\\
If $p(z)=\sum_j c_jz^j$, then $p(A)=\sum\limits_j c_jA^j =S\left(\sum\limits_j c_j\Lambda^j\right)S^{-1} =S p(\Lambda) S^{-1}$.\\
Since $\Lambda$ is diagonal, $p(\Lambda)=\operatorname{diag}(p(\lambda_1),\ldots,p(\lambda_n))$.\\
\hlt{Structural Commentary}: Diagonalization reduces $p(A)$ to scalar evaluation on eigenvalues.\\
\hlt{Key Insight}: Compute functions of $A$ by applying the function to the eigenvalues.
\end{proof}
\end{theorem}
